% Table S5 — EPMN native vs shared protocol
\begin{table}[t]
\centering
\setlength{\tabcolsep}{6pt}
\caption{%
EPMN native-recipe versus shared-protocol comparison (mean $\pm$ standard deviation across subjects, LOSO) at the three-channel montage. The native recipe follows the training procedure of Wei et al.\ \cite{wei2022}: input resampled to 64 time points; SGD optimizer (learning rate $10^{-4}$, momentum 0.9) with learning-rate halving every 5 epochs; episodic training with $N_q = 12$ query trials and $N_s = 18$ support trials per class, re-sampled each episode; majority-class down-sampling with unweighted cross-entropy; no augmentation. The shared protocol uses the AdamW optimizer, cosine-annealed learning rate, temporal jitter, and additive noise described in Section~2.5 of the main text. Both protocols use the same two-phase subject-level validation scaffolding and LOSO folds. Bold marks the higher mean per dataset. $\Delta$ = shared-protocol minus native AUROC (positive indicates the shared protocol is higher). The comparison was run on the three datasets for which the native EPMN episodic recipe could be implemented from the original description. EPMN: ERP Prototypical Matching Net; abbreviations otherwise as in Table~S1.%
}
\label{sup:tab-epmn-native}
\begin{tabular}{lccc}
\toprule
Dataset & Native AUROC & Shared-protocol AUROC & $\Delta$ \\
\midrule
ERN & 0.776 $\pm$ 0.121 & \textbf{0.795 $\pm$ 0.096} & \textit{+0.019} \\
P300 & 0.629 $\pm$ 0.078 & \textbf{0.674 $\pm$ 0.077} & \textit{+0.045} \\
N400 & 0.565 $\pm$ 0.076 & \textbf{0.573 $\pm$ 0.077} & \textit{+0.008} \\
\midrule
\textit{Mean} & 0.656 & \textbf{0.681} & \textbf{+0.024} \\
\bottomrule
\end{tabular}
\end{table}
